\documentclass[11pt,a4paper]{article}
\usepackage[T1]{fontenc}
\usepackage{amsmath,amssymb,amsthm}
\usepackage[margin=28mm]{geometry}
\usepackage[hidelinks]{hyperref}
\newtheorem{theorem}{Theorem}
\newtheorem{lemma}[theorem]{Lemma}
\newtheorem{corollary}[theorem]{Corollary}
\newtheorem{proposition}[theorem]{Proposition}
\theoremstyle{remark}
\newtheorem{remark}[theorem]{Remark}
\newcommand{\Z}{\mathbb Z}
\newcommand{\R}{\mathbb R}
\newcommand{\Nzero}{\Z_{\geq0}}
\newcommand{\card}[1]{\lvert#1\rvert}
\newcommand{\norm}[1]{\lVert#1\rVert_\infty}
\newcommand{\floor}[1]{\lfloor#1\rfloor}
\newcommand{\ceil}[1]{\lceil#1\rceil}
\newcommand{\dd}{\delta}
\title{Cumulative growth of corner cuts\\in fixed dimension}
\author{}
\date{17 October 2026}
\begin{document}
\maketitle
\begin{center}
\small Working mathematical draft. The proof has undergone an internal audit;
the priority assessment and submission review remain incomplete.
\end{center}

\begin{abstract}
Let \(C_d(n)\) count the \(n\)-point subsets of the nonnegative integer
lattice that can be strictly separated from their complements by an affine
hyperplane. For every fixed \(d\geq2\), we prove
\[
 \sum_{j=1}^{n}C_d(j)=\Theta_d\!\bigl(n^d(\log n)^{d-1}\bigr).
\]
The upper bound follows from Wagner's estimate. His lifting lemma then
gives \(C_d(n)=\Omega_d(n^{d-1}(\log n)^{d-2})\) for \(d\geq4\),
improving the lower logarithmic exponent in that published baseline.
The construction counts primitive lattice simplices in separated boxes,
converts their affine spans to vertices of a weight-space arrangement,
and sums over disjoint aspect ranges. A finite divisibility argument
provides uniform constants without an interchange of density limits
and infinite unions. The exact-size three-dimensional gap is not resolved.
\end{abstract}

\section{The counting problem and its scope}
A nonempty finite set \(D\subset\Nzero^d\) is a \emph{corner cut} if
\[
 D=\{x\in\Nzero^d:w\cdot x<t\}
 \quad\text{for some }w\in\R_{>0}^d,\ t>0,
\]
with no lattice point on the cutting hyperplane. Write \(C_d(n)\) for the
number of such sets of cardinality \(n\), and put
\[
 A_d(n)=\sum_{j=1}^{n}C_d(j).
\]
We count sets, rather than the orders of their elements or the weights
representing them. Every cut is an order ideal for the coordinate order.
In particular, every coordinate of an element of an \(n\)-point cut
is at most \(n-1\), so these counts are finite.

These objects occur as the exponent sets of initial segments of smooth
numbers. For distinct primes \(p_1,\ldots,p_d\), comparing the products
\(\prod_i p_i^{x_i}\) is equivalent to comparing
\(\sum_i x_i\log p_i\). An exponent prefix is consequently a corner cut.
Our theorem concerns the full geometric class; it does not give a
prime-distribution estimate or a faster random-access algorithm for a
fixed prime set.

Corner cuts were introduced in connection with state polyhedra and
Gr\"obner bases by Onn and Sturmfels \cite{OS}. The planar order
\(C_2(n)=\Theta(n\log n)\) is due to Corteel, R\'emond, Schaeffer and
Thomas \cite{CRST}. For fixed \(d\), Wagner \cite{W02} proved
\[
 \Omega_d(n^{d-1}\log n)\leq C_d(n)
 \leq O_d(n^{d-1}(\log n)^{d-1}).
 \tag{1}\label{eq:baseline}
\]
His thesis \cite[Chapter 3]{W03} gives the lower and upper arguments
and suggests seeking greater completion multiplicity. Our new step is
a lower bound for bounded cardinality, obtained by counting independent
aspect ranges. Lifting is used only after this step.

\begin{theorem}\label{thm:main}
For each fixed integer \(d\geq2\),
\[
 A_d(n)=\Theta_d\!\bigl(n^d(\log n)^{d-1}\bigr).
\]
More explicitly, let
\[
 R_0(d)=4d^2(d+1),\qquad n_0(d)=2^d4^{R_0(d)}
\]
and
\[
 \kappa_d=
 \frac{1}{
 8^{d-1}(100d)^{d^2}(d-1)![4(d+1)]^{d-1}2^{d(d+2)}}.
\]
Then, for every integer \(n\geq n_0(d)\),
\[
 A_d(n)\geq\kappa_d n^d(\log_4 n)^{d-1}.
 \tag{2}\label{eq:explicit}
\]
\end{theorem}

\begin{corollary}\label{cor:exact}
For each fixed \(d\geq4\),
\[
 C_d(n)=\Omega_d\!\bigl(n^{d-1}(\log n)^{d-2}\bigr).
\]
\end{corollary}
This leaves one logarithmic factor between the lower bound and
\eqref{eq:baseline}. In dimension three, the corollary recovers only
the earlier lower order \(n^2\log n\). The cumulative theorem also
precludes a uniform reduction of the logarithmic exponent in the upper
bound; its positive-proportion consequence is stated in Section~\ref{sec:consequences}.

\section{Primitive simplices in axis boxes}
Fix an ambient integer \(d\geq2\) and set \(\dd=1/(100d)\).
For integer side lengths \(L_1,\ldots,L_d\), put
\[
 m=\min_i L_i,\quad M=\max_i L_i,\quad Q=\prod_{i=1}^{d}L_i.
\]
Define \(B_i\subset\Nzero^d\) by taking the \(j\)th coordinate interval to be
\[
 [L_j,L_j+\floor{\dd L_j}]\quad(j=i),\qquad
 [0,\floor{\dd L_j}]\quad(j\ne i).
\]
A \(d\)-by-\((d-1)\) integer matrix is called primitive if the greatest
common divisor of its maximal minors is one.

\begin{lemma}\label{lem:primitive}
Assume \(M\leq m^2\) and \(m\geq2^{8d^2}\).
At least a fraction \(8^{-(d-1)}\) of the ordered choices
\(p_i\in B_i\) have primitive difference matrix
\[
 V=[p_2-p_1\ \cdots\ p_d-p_1].
\]
\end{lemma}
The minor criterion and its equivalence with full column rank modulo
every prime are standard \cite[Propositions 3--4]{MW}.
The point requiring proof here is uniformity for these moving,
possibly anisotropic boxes.

\begin{proof}
Keep the same \(\dd\) throughout a dimension induction.
Let \(E_r\) be the event that the first \(r\) points, projected onto the
first \(r\) coordinates, have primitive difference matrix.
The event \(E_1\) is automatic, with empty minor equal to one.
At step \(2\leq r\leq d\), fix a primitive configuration of the previous
\(r-1\) projected points. Its difference matrix \(U\) has size
\((r-1)\)-by-\((r-2)\). Its vector \(a_0\) of signed maximal minors
is primitive, and \(U\) has rank \(r-2\) modulo every prime.

Let \(u\) be the first \(r-1\) coordinates of \(p_r-p_1\), and put
\(D=\det[U\ u]\). This determinant is nonzero for every allowed choice.
Indeed, the normalized row matrix of the previous \(r-1\) projected
points differs from the identity by at most \((r-1)\dd\) in maximum
row-sum norm. The normal of their affine span, normalized to height one,
has scaled coordinates within
\((r-1)\dd/(1-(r-1)\dd)\) of one and is positive.
The new projected point has every scaled coordinate at most \(\dd\);
it is strictly on the origin side of that affine span.
Also,
\[
 |D|\leq (r-1)!\,2^{r-1}M^{r-1}
       \leq2^{d^2}m^{2(d-1)}.                    \tag{3}\label{eq:D}
\]

Every coordinate interval contains more than \(\dd m\) integers.
For a prime \(\ell\), the equation \(D=0\pmod\ell\) is a nontrivial
linear equation in the new projected point: at least one coordinate
of \(a_0\) is nonzero modulo \(\ell\).
Its probability is at most \(1/\ell+a/m\), where \(a=100d\).

Suppose \(\ell\mid D\). Since \(U\) has full column rank modulo
\(\ell\), there is a unique vector \(\lambda\) with
\(u=U\lambda\pmod\ell\). Write \(z\) for the \(r\)th row of the old
columns and \(u_r\) for the \(r\)th coordinate of the new column.
The full \(r\)-by-\((r-1)\) matrix loses rank modulo \(\ell\) precisely if
\[
 u_r=z\lambda\pmod\ell.                         \tag{4}\label{eq:residue}
\]
Conditional on everything except the last coordinate of \(p_r\),
this is one residue class, of probability at most \(1/\ell+a/m\).
Thus the probability of failure at any particular prime is at most
\((1/\ell+a/m)^2\).

Put \(s=\floor{(\log_2 m)/2}\), \(P=2^s\).
The small-prime union bound gives
\begin{align}
 \Pr(\text{failure at a prime }\ell\leq P)
 &\leq\sum_{\substack{\ell\leq P\\\ell\ {\rm prime}}}
                 (1/\ell+a/m)^2 \notag\\
 &\leq\frac34+\frac{2a(s+1)}{2^{2s}}+
                      \frac{a^2}{2^{3s}}.       \tag{5}\label{eq:small}
\end{align}
Here \(\sum_{j=2}^{\infty}j^{-2}<3/4\) follows by comparison with
\(1/4+\int_2^\infty x^{-2}\,dx\).
The cross term uses \(H_P\leq1+\log P\leq1+s\), and the constant
term uses at most \(P\) primes.

For large primes, first fix the new projected coordinates.
Every bad prime \(\ell>P\) divides the same nonzero integer \(D\).
By \eqref{eq:D}, their number is at most
\[
 \frac{\log_2|D|}{s}
 \leq4(d-1)+\frac{d^2+4(d-1)}{s}
 \leq b:=4d-3
 \quad\text{if }s\geq d^2+4(d-1).
\]
For each of these finitely many primes, \eqref{eq:residue} bounds the
conditional last-coordinate failure probability by \(2^{-s}+a2^{-2s}\).
Their total contribution is therefore at most \(b(2^{-s}+a2^{-2s})\).

Define
\[
 E_d(s)=\frac{2a(s+1)}{2^{2s}}+\frac{a^2}{2^{3s}}
                  +b(2^{-s}+a2^{-2s}).          \tag{6}\label{eq:error}
\]
This decreases for \(s\geq1\).
For every integer \(d\geq2\), the choice \(S_d=4d^2\) satisfies
\(S_d\geq d^2+4(d-1)\) and \(E_d(S_d)<2^{-10}<1/8\).
For an explicit verification valid in every dimension, use
\(a\leq2^{4d}\), \(b\leq2^{2d}\) and \(4d^2+1\leq2^{3d}\).
They give
\[
 E_d(4d^2)\leq
 2^{1+7d-8d^2}+2^{8d-12d^2}
       +2^{2d-4d^2}+2^{6d-8d^2}<2^{-10}.
\]
The four exponents decrease for \(d\geq2\), and at \(d=2\)
are \(-17,-32,-12,-20\).

Consequently the conditional probability of a nonprimitive extension
is at most \(3/4+E_d(s)\leq7/8\).
To justify the conditioning across steps, the events
\(E_1,\ldots,E_{r-1}\) depend only on previous projected coordinates.
After these are fixed, the old \(r\)th coordinates and the new point
retain their product distributions. Thus
\[
 \Pr(E_r\mid E_1,\ldots,E_{r-1})\geq1/8.
\]
Multiplication proves that \(\Pr(E_1\cap\cdots\cap E_d)\geq8^{-(d-1)}\).
This intersection implies \(E_d\), as required.
All bounds on large primes were finite before any limit was taken.
\end{proof}

\section{From simplices to cuts}
We first recall a useful elementary fact about arrangements, including
the case of dependent hyperplanes.
\begin{lemma}\label{lem:chambers}
For a finite affine-hyperplane arrangement in a nonempty open convex
domain in \(\R^k\), the number of chambers is at least the number
of interior vertices.
\end{lemma}
\begin{proof}
In dimension one, \(v\) distinct interior points give \(v+1\) chambers.
For the induction step, discard duplicate hyperplanes and insert the
remaining ones in a fixed order. Inserting \(H\) creates exactly the
number of chambers of the restricted arrangement in \(H\) intersected
with the domain. This follows because every old chamber is open and
convex and is split exactly when it meets \(H\).

Assign each vertex to its last incident hyperplane \(H\).
The incident normals span \(\R^k\), and the nonzero normal of \(H\)
extends to a basis among them. Hence the assigned vertex is also an
interior vertex of the restriction to \(H\).
The induction hypothesis bounds the number assigned to \(H\) by the
number of new chambers. Summing gives the result.
\end{proof}

\begin{lemma}\label{lem:local}
Under the hypotheses of Lemma~\ref{lem:primitive}, at least
\[
 \rho_d Q^d,\qquad
 \rho_d=8^{-(d-1)}(100d)^{-d^2},                 \tag{7}\label{eq:local}
\]
distinct corner cuts are represented by weights in
\[
 \Omega_L=\prod_{i=1}^{d}\left(\frac9{10L_i},\frac{11}{10L_i}\right),
\]
using threshold one. Each has at most \((3/2)^d Q\) points.
Its maximum \(i\)th coordinate on the \(i\)th axis lies in
\((4L_i/5,6L_i/5)\).
\end{lemma}
\begin{proof}
Let \(p_i\in B_i\) have primitive difference matrix \(V\).
Its primitive cofactor normal is \(a\).
There is an integer vector \(z\) with \(a\cdot z=1\);
the matrix \([V\ z]\) is unimodular, after choosing the cofactor
orientation. The columns of \(V\) therefore form a basis of
\(\{x\in\Z^d:a\cdot x=0\}\).

After row scaling, the submatrix of \(V\) in coordinates \(2,\ldots,d\)
differs from the identity by at most \((d-1)\dd\) in maximum row-sum
norm. If two integer points in one \(B_i\) belong to the selected
affine span, their difference is \(Vc\) for an integer vector \(c\).
The scaled coordinate difference is at most \(\dd\), so
\[
 \norm{c}\leq\frac{\dd}{1-(d-1)\dd}<1.
\]
Thus \(c=0\). The plane meets each box in exactly its selected point,
so distinct primitive choices give distinct planes.

The normalized row matrix of all \(d\) points differs from the
identity by at most \(d\dd\).
The unique solution of \(p_i\cdot w=1\) satisfies
\[
 |L_iw_i-1|\leq\frac{d\dd}{1-d\dd}<1/10.
\]
In particular it is positive and lies in \(\Omega_L\).
Each selected plane gives a distinct interior vertex of the arrangement
\(x\cdot w=1\).
This arrangement can be taken finite: only nonnegative integer points
with \(x_i\leq\ceil{10L_i/9}\) for every \(i\) can meet \(\Omega_L\);
all other points have weight greater than one throughout that domain.

Every chamber gives a corner cut. Different chambers give different
sets, since a fixed collection of strict signs has convex intersection
with \(\Omega_L\). By Lemma~\ref{lem:chambers}, the number of these cuts
is at least the number of selected vertices.
Each \(B_i\) has at least \(\dd^d Q\) points, so
Lemma~\ref{lem:primitive} proves \eqref{eq:local}.

A coordinate of an included point is smaller than \(10L_i/9\).
There are fewer than \(10L_i/9+1\leq(3/2)L_i\) possible integer values
when \(L_i\geq3\). This proves the cardinality bound.
For a chamber weight, the maximum axis coordinate is \(\floor{1/w_i}\);
the chamber condition excludes equality \(1/w_i\in\Z\).
It lies strictly between \(4L_i/5\) and \(6L_i/5\) when \(L_i\geq10\).
Both side conditions follow from the stated threshold.
\end{proof}

\section{Summing aspect ranges}
\begin{proof}[Proof of Theorem~\ref{thm:main}]
The upper bound follows by summing \eqref{eq:baseline}.
For the lower bound, take \(R\) maximal with
\[
 Q=4^R\leq n/2^d.
\]
Choose any integer tuple with
\[
 r_1+\cdots+r_d=R,\qquad r_i\geq\ceil{R/(d+1)},
 \quad\text{and put } L_i=4^{r_i}.
\]
Then
\[
 Q^{1/(d+1)}\leq L_i\leq Q^{2/(d+1)}.
\]
In particular \(\max L_i\leq(\min L_i)^2\).
When \(n\geq n_0(d)\), we have \(R\geq R_0(d)\) and
\(\min L_i\geq2^{8d^2}\), so Lemma~\ref{lem:local} applies.

All the resulting cuts have at most
\((3/2)^dQ<n\) points. No cut is counted for two distinct scale tuples:
such tuples differ by a factor of at least four in some \(L_i\),
whereas the observable maximum axis coordinate belongs to
\((4L_i/5,6L_i/5)\). These ranges are disjoint.

Write \(h=R-d\ceil{R/(d+1)}\).
There are \(\binom{h+d-1}{d-1}\) tuples.
For \(R\geq R_0(d)\),
\[
 h\geq R/(d+1)-d\geq R/[2(d+1)],
\]
and hence
\[
 \binom{h+d-1}{d-1}
 \geq\frac{R^{d-1}}{[2(d+1)]^{d-1}(d-1)!}.
\]
Also \(Q>n/2^{d+2}\) and
\(R\geq\log_4 n-d/2-1\geq(\log_4 n)/2\).
Multiply this tuple count by \(\rho_d Q^d\) from
\eqref{eq:local} to obtain \eqref{eq:explicit}.
\end{proof}

For clarity, we give the lifting argument that transfers this result to
a prescribed size in one more dimension. The argument is due to
Wagner \cite[Lemma 3.5]{W03}.
\begin{lemma}\label{lem:lifting}
For \(d\geq1\) and \(n\geq1\),
\[
 C_{d+1}(n)\geq\sum_{j=1}^{n}C_d(j).
\]
\end{lemma}
\begin{proof}
Take a \(j\)-point corner cut \(T\) in the plane \(x_{d+1}=0\), and
choose a strict representation \(w\cdot x<1\).
Within its open representation set we can choose the coordinates
of \(w\), together with one, to be linearly independent over \(\mathbb Q\).
For \(u>0\) consider
\[
 D_u=\{(x,z)\in\Nzero^{d+1}:w\cdot x+uz<1\}.
\]
Its intersection with \(z=0\) is always \(T\).
For \(u>1\), \(D_u=T\).
As \(u\) decreases to zero, points enter at the values
\((1-w\cdot x)/z\), with \(x\in T\), \(z\geq1\).
These entry values are distinct: equality for two pairs implies a
rational linear relation among \(1,w_1,\ldots,w_d\), and then forces
both \(z\) and \(x\) to agree.
There are only finitely many such values above any fixed positive
bound, and the cardinality tends to infinity. Thus some open interval
of \(u\)'s gives cardinality exactly \(n\), with no boundary point.
Cuts with distinct sections \(T\) are different.
\end{proof}
\begin{proof}[Proof of Corollary~\ref{cor:exact}]
Apply Lemma~\ref{lem:lifting} to \(A_{d-1}(n)\), and then use
Theorem~\ref{thm:main}.
\end{proof}

\section{Consequences and limitations}\label{sec:consequences}
\begin{corollary}
For every fixed \(d\geq2\), there are \(c_d,\eta_d>0\) such that, for
every sufficiently large \(N\), at least \(\eta_dN\) integers \(j\leq N\)
satisfy
\[
 C_d(j)\geq c_d j^{d-1}(\log j)^{d-1}.
\]
\end{corollary}
\begin{proof}
Choose \(a,b>0\) with \(A_d(N)\geq aN^d(\log N)^{d-1}\) for large \(N\)
and \(C_d(j)\leq bj^{d-1}(\log j)^{d-1}\) for large \(j\).
The finitely many exceptional terms contribute \(O_d(1)\).
For \(0<c<a/4\), the sum over indices failing the claimed threshold is
at most \(cN^d(\log N)^{d-1}+O_d(1)\).
Every other term is at most \(bN^{d-1}(\log N)^{d-1}\).
At least a fixed positive multiple of \(N\) terms are therefore required.
\end{proof}

There is an established geometric interpretation.
For a cut \(D\), put \(\alpha_D=\sum_{x\in D}x\).
For fixed size \(n\), these vectors are distinct and each is a vertex
of their convex hull. Indeed, weights strictly cutting off \(D\)
make \(w\cdot\alpha_D\) smaller than \(w\cdot\alpha_E\) for every
other \(n\)-point set \(E\): each replacement point has greater
weight than every removed point. Thus \(\alpha_D\) is uniquely exposed.
Conversely every vertex of the finite hull is one of its generators.
The vertex count of the corner-cut polytope is consequently \(C_d(n)\).
This correspondence is known \cite{OS,BMORW}; our contribution is the
quantitative lower count, not the dictionary.
For generic experimental designs the same count describes their
corner-cut models \cite[Definition 7 and Theorem 2]{BMORW}.

The cumulative result gives sharp-order growth on average and the
improved exact-size bound in dimensions at least four.
It supplies no asymptotic leading constant and no pointwise
\(C_3(n)\gg n^2(\log n)^2\) statement.
It also gives no new enumeration algorithm, no uniform claim as
\(d\) grows with \(n\), and no fixed-prime unranking complexity improvement.

\section{Prior results and reproducibility}
The following distinctions delimit the contribution being assessed.
The unrestricted density of rectangular primitive integer matrices is
known \cite{MRW,MW}. The later paper repairs a defect in the earlier
density proof. A centered-cube density alone does not assert the
uniform estimate for our moving boxes; Lemma~\ref{lem:primitive} gives
a direct finite proof. We claim no new minor criterion or unrestricted
density formula.

Chen, Feng, Liu and Wu \cite[Theorem 1]{CFLW} give finite probability
bounds for appending independent rows to a prescribed primitive matrix.
Their random entries use a common interval of length at least the
maximum prescribed entry. Our boxes have translated intervals of
different lengths, with their shortest length potentially much smaller
than the largest entry. That theorem is not a direct substitute for
Lemma~\ref{lem:primitive}. Prime-by-prime rank analysis and finite
determinant cutoffs are established methods; our claimed contribution
is their use in the multiscale corner-cut count.

Finite-grid threshold encodings and planar counts are related
background \cite{GPZ,ZZ}. Their parameters are the ambient grid and,
in the latter case, intersections of two halfplanes. The argument here
additionally controls the number of included points, restricts the
positive intercepts, and sums disjoint aspect ranges.
A difference in definitions alone is not evidence of novelty.

The Hilbert zonotope refines state polyhedra and supports polynomial-time
computation of universal Gr\"obner bases in fixed dimension
\cite[Proposition 3.2 and Theorem 3.5]{BOT}.
Its count is an upper bound for a common refinement, not the lower
bounded-cardinality count proved here.
Likewise the lattice-subspace covering bounds of
Balko, Cibulka and Valtr \cite[Theorems 3--5]{BCV} concern another
counting parameter. Neither result is being recast as a new theorem.

Elizalde and Galv\'an \cite[Section 8.2]{EG} discuss higher-dimensional
pyramidal partitions and refer to the corner-cut growth literature.
For this draft, the relevant results in their 2025 journal article are
provisionally assumed unchanged from arXiv:2312.16353v1, as authorized
for the project. The final-version comparison is recorded separately
for later checking. This assumption does not establish priority.

\paragraph{Current priority status.}
The proof establishes the stated inequalities and their strict
improvement over the inspected lower bound in \eqref{eq:baseline}.
The literature audit has not established an unrestricted claim of
priority or submission readiness. In particular, the complete 2002
journal text has not been inspected; the proofs were compared with
the author's full thesis chapter. The 2007 grid-encoding paper was
compared at the abstract level only. The detailed audit accompanies
the source at \path{results/research/corner_cut_growth_audit.json}.
This paragraph must be resolved before the draft is described as
submission-ready.

\paragraph{Exact checks.}
The standard-library Python checker
\path{scripts/run_corner_cut_growth_checks.py}
uses independent determinants, modular Gaussian elimination and exact
rational feasibility. It checks primitive-plane injectivity, the
conditional residue bound, a nonprimitive-projection negative control,
and a degenerate arrangement with 18 interior vertices and 48 chambers.
It exhausts stated small boxes in dimensions two through four and uses
deterministic samples in dimensions five and six.
The explicit constants and 352 scale-rounding cases in dimensions
two through twelve are checked with rational arithmetic.
These computations can falsify steps; they do not prove the uniform
theorem or its novelty.

Run, from the repository root,
\begin{verbatim}
python3 -B scripts/run_corner_cut_growth_checks.py
\end{verbatim}
The report records fixtures, witnesses, the exact command and SHA-256
hashes of its sources in
\path{results/research/corner_cut_growth_checks.json}.
The proof above is independent of the computation.
The other project kernels and public rank/unrank interface are unchanged.

\begin{thebibliography}{99}
\raggedright
\bibitem{OS}
S.~Onn and B.~Sturmfels,
Cutting corners,
\emph{Advances in Applied Mathematics} 23 (1999), 29--48.
\href{https://doi.org/10.1006/aama.1999.0645}{doi:10.1006/aama.1999.0645}.

\bibitem{CRST}
S.~Corteel, G.~R\'emond, G.~Schaeffer and H.~Thomas,
The number of plane corner cuts,
\emph{Advances in Applied Mathematics} 23 (1999), 49--53.
\href{https://doi.org/10.1006/aama.1999.0646}{doi:10.1006/aama.1999.0646}.

\bibitem{W02}
U.~Wagner,
On the number of corner cuts,
\emph{Advances in Applied Mathematics} 29 (2002), 152--161.
\href{https://doi.org/10.1016/S0196-8858(02)00014-3}{doi:10.1016/S0196-8858(02)00014-3}.

\bibitem{W03}
U.~Wagner,
\emph{On \(k\)-Sets and Their Applications},
doctoral thesis, ETH Z\"urich (2003), Chapter 3.
\href{https://doi.org/10.3929/ethz-a-004708408}{doi:10.3929/ethz-a-004708408}.

\bibitem{MRW}
G.~Maze, J.~Rosenthal and U.~Wagner,
Natural density of rectangular unimodular integer matrices,
\emph{Linear Algebra and its Applications} 434 (2011), 1319--1324.
\href{https://doi.org/10.1016/j.laa.2010.11.015}{doi:10.1016/j.laa.2010.11.015}.

\bibitem{MW}
G.~Micheli and V.~Weger,
On rectangular unimodular matrices over the algebraic integers,
\emph{SIAM Journal on Discrete Mathematics} 33 (2019), 425--437.
\href{https://doi.org/10.1137/18M1177093}{doi:10.1137/18M1177093};
\href{https://arxiv.org/abs/1803.08785}{arXiv:1803.08785}.

\bibitem{BMORW}
Y.~Berstein, H.~Maruri-Aguilar, S.~Onn, E.~Riccomagno and H.~Wynn,
Minimal average degree aberration and the state polytope for experimental designs,
\emph{Annals of the Institute of Statistical Mathematics} 62 (2010), 673--698.
\href{https://doi.org/10.1007/s10463-010-0291-8}{doi:10.1007/s10463-010-0291-8}.

\bibitem{BOT}
E.~Babson, S.~Onn and R.~R.~Thomas,
The Hilbert zonotope and a polynomial time algorithm for universal Gr\"obner bases,
\emph{Advances in Applied Mathematics} 30 (2003), 529--544.
\href{https://doi.org/10.1016/S0196-8858(02)00509-2}{doi:10.1016/S0196-8858(02)00509-2};
\href{https://arxiv.org/abs/math/0207135}{arXiv:math/0207135}.

\bibitem{GPZ}
S.~Ghilezan, J.~Pantovi\'c and J.~\v{Z}uni\'c,
Separating points by parallel hyperplanes---characterization problem,
\emph{IEEE Transactions on Neural Networks} 18 (2007), 1356--1363.
\href{https://doi.org/10.1109/TNN.2007.891678}{doi:10.1109/TNN.2007.891678}.

\bibitem{ZZ}
E.~Zamaraeva and J.~\v{Z}uni\'c,
A characterization of 2-threshold functions via pairs of prime segments,
\emph{Theoretical Computer Science} 919 (2022), 1--17.
\href{https://doi.org/10.1016/j.tcs.2022.03.025}{doi:10.1016/j.tcs.2022.03.025};
\href{https://arxiv.org/abs/2007.03986}{arXiv:2007.03986}.

\bibitem{BCV}
M.~Balko, J.~Cibulka and P.~Valtr,
Covering lattice points by subspaces and counting point-hyperplane incidences,
\href{https://arxiv.org/abs/1703.04767}{arXiv:1703.04767}.

\bibitem{CFLW}
J.~Chen, Y.~Feng, Y.~Liu and W.~Wu,
On the probability of generating a primitive matrix,
\emph{Journal of Systems Science and Complexity} (2024).
\href{https://doi.org/10.1007/s11424-024-1313-6}{doi:10.1007/s11424-024-1313-6}.
Inspected version:
\href{https://arxiv.org/abs/2105.05383v2}{arXiv:2105.05383v2}.

\bibitem{EG}
S.~Elizalde and A.~B.~Galv\'an,
Triangular partitions: enumeration, structure, and generation,
\emph{Advances in Applied Mathematics} 163 (2025), 102807.
\href{https://doi.org/10.1016/j.aam.2024.102807}{doi:10.1016/j.aam.2024.102807}.
Inspected version:
\href{https://arxiv.org/abs/2312.16353v1}{arXiv:2312.16353v1}.
\end{thebibliography}
\end{document}
